\chapter{Аудит безаномального носителя асимптотически-свободной моды}
\label{ch:v10-discrete-resolution-af-carrier-audit}

Текущие неабелевы сектора не дают требуемого коэффициента:
\begin{equation}
 b_{SU(2)_{\rm SM}}=-\frac{19}{6},\qquad
 b_{SU(3)_{\rm SM}}=-7,
\end{equation}
тогда как полюсная трансмутация требует $b=-2$. Поэтому наличие
асимптотической свободы ещё не означает наличие нужного носителя.

В ограниченном классе $SU(N)$ с фундаментальными дираковскими фермионами
и комплексными скалярами минимальные неотрицательные решения уравнения
\begin{equation}
 4n_D+n_s=22N-12
\end{equation}
для $N=2,\ldots,6$ равны
\begin{equation}
 \begin{array}{c|cc|c}
 N&n_D&n_s&N+n_D+n_s\\\hline
 2&8&0&10\\3&13&2&18\\4&19&0&23\\5&24&2&31\\6&30&0&36
 \end{array}.
\end{equation}
Следовательно, в этом классе единственный минимальный кандидат ---
$SU(2)$ с восемью дираковскими фундаменталами. Он содержит $16$
вейлевских дублетов, поэтому проходит глобальный $SU(2)$ anomaly parity.

Однако минимальность не является вложением. Кандидат отсутствует в
типизированном K43-носителе, не имеет морфизма к ранговому проектору
трансмутированной моды и не выбран родителем кратностей. Аналогично,
имеющиеся SM-$SU(2)$ и SM-$SU(3)$ внутренни и безаномальны, но имеют
неправильные beta-коэффициенты.

Проверка одиннадцати кандидатов по шести критериям даёт матрицу ранга $6$
и оценки
\begin{equation}
 (4,4,4,3,3,3,3,2,3,3,5).
\end{equation}
Полных проходов нет. Формальный AF-portal получает $5/6$, но подставляет
новый сектор без родительского происхождения.

\begin{theorem}[Минимальный условный носитель определён]
В фундаментальном $SU(N)$-классе минимальная безаномальная реализация
$b=-2$ есть $SU(2)$ с восемью дираковскими дублетами. Это условная
архитектура, а не уже существующий сектор теории.
\end{theorem}

\begin{nogocriterion}[Минимальность не заменяет типизированное вложение]
Кандидат нельзя считать физическим, пока его $SU(2)$-действие, шестнадцать
вейлевских компонент, калибровочный оператор и связь с полюсной модой не
реализованы внутри общего K43-носителя и общего родителя.
\end{nogocriterion}

\begin{protocol}\begin{enumerate}
\item условная минимальность: \textbf{$1/1$};
\item глобальная anomaly parity: \textbf{$1/1$};
\item полные проходы кандидатов: \textbf{$0/11$};
\item физическое происхождение: \textbf{$0/3$};
\item следующий гейт: \textbf{K43-вложение $SU(2)+8D$}.
\end{enumerate}\end{protocol}

Машинный сертификат сохранён в
\path{s2t_v10_cell_birth_four_volume_nonequilibrium_bath_discrete_resolution_transmuted_mode_asymptotically_free_anomaly_free_carrier_candidate_audit_gate_results.json}.